\section{Mathematical Background}
\label{sec:background}

\subsection{The Lie--Trotter product formula}
\label{sec:bg-trotter}

Let $\A$ be a unital Banach algebra over $\mathbb{K} \in \{\mathbb{R}, \mathbb{C}\}$
with norm $\norm{\cdot}$ satisfying $\norm{1} = 1$ and $\norm{ab} \le \norm{a}\norm{b}$.
The exponential
\[
  e^a = \sum_{k=0}^{\infty} \frac{a^k}{k!}
\]
converges absolutely for every $a \in \A$.

Trotter~\cite{trotter1959product} proved that for bounded operators $A, B$,
\[
  e^{A+B} = \lim_{n\to\infty} \bigl(e^{A/n}\,e^{B/n}\bigr)^n.
\]
The standard proof proceeds in two steps:
\begin{enumerate}
  \item \emph{Telescoping}: writing $M = \max\{\norm{X},\norm{Y}\}$,
    the identity
    $X^n - Y^n = \sum_{k=0}^{n-1} X^k(X-Y)Y^{n-1-k}$ gives
    $\norm{X^n - Y^n} \le n\,\norm{X-Y}\,M^{n-1}$.

  \item \emph{Step error}: the single-step approximation satisfies
    $\norm{e^a\,e^b - e^{a+b}} = O(\norm{a}\,\norm{b})$.
\end{enumerate}
Setting $X = e^{A/n}\,e^{B/n}$ and $Y = e^{(A+B)/n}$, the $O(1/n^2)$ step
error is amplified by the factor $n$ from telescoping, yielding $O(1/n)$
overall convergence.

\subsection{Higher-order splitting methods}
\label{sec:bg-splitting}

The symmetric \emph{Strang splitting}~\cite{strang1968construction} uses
\[
  S_2(t) = e^{tA/2}\,e^{tB}\,e^{tA/2}
\]
to achieve $O(t^3)$ local error and $O(1/n^2)$ global convergence.
Suzuki~\cite{suzukiGeneralTheoryFractal1991} showed that higher-order integrators can be
constructed recursively: the fourth-order formula
\[
  S_4(t) = S_2(pt)^2\,S_2\bigl((1-4p)t\bigr)\,S_2(pt)^2
  \qquad \text{with } p = \frac{1}{4 - 4^{1/3}}
\]
achieves $O(t^5)$ local error.

Both formulas generalize to multiple operators $A_1, \ldots, A_m$.
The first-order product $(e^{A_1/n} \cdots e^{A_m/n})^n$ converges at rate
$O(1/n)$, while the palindromic product
\[
  S_2^{(m)}(t) = e^{tA_1/2}\,e^{tA_2/2}\cdots e^{tA_m}\cdots e^{tA_2/2}\,e^{tA_1/2}
\]
achieves $O(1/n^2)$.

\subsection{Commutator-scaling error bounds}
\label{sec:bg-commutator}

The standard step error $\norm{e^a\,e^b - e^{a+b}} \le 2\norm{a}\norm{b}e^{\norm{a}+\norm{b}}$
treats $A$ and $B$ as unrelated. When $A$ and $B$ nearly commute---as is
typical for spatially local Hamiltonians---the error is much smaller.

Childs et al.~\cite{childsTheoryTrotterError2021} showed that the leading error term is
controlled by the commutator $\comm{B}{A} = BA - AB$ rather than the product
$\norm{A}\norm{B}$. At first order,
\begin{equation}\label{eq:comm-scaling-bg}
  \norm{e^{tB}\,e^{tA} - e^{t(A+B)}}
    \le \frac{\norm{\comm{B}{A}}}{2}\,t^2\,e^{t(\norm{A}+3\norm{B})}.
\end{equation}
For the Strang splitting, the leading commutator $\comm{B}{A}$ cancels by
symmetry, and the error is controlled by double commutators:
\begin{equation}\label{eq:strang-comm-bg}
  \norm{S_2(t) - e^{tH}}
    \le \Bigl(\frac{\norm{\comm{B}{\comm{B}{A}}}}{12}
      + \frac{\norm{\comm{A}{\comm{A}{B}}}}{24}\Bigr)\,t^3.
\end{equation}

These commutator-scaling bounds are physically significant: for a lattice
Hamiltonian $H = \sum_j h_j$ where each $h_j$ acts on a few neighboring sites,
$\norm{\comm{h_j}{h_k}} = 0$ whenever the supports of $h_j$ and $h_k$ are
disjoint, so the commutator-based bound can be polynomially tighter than the
naive product bound~\cite{childsTheoryTrotterError2021}.

\subsection{Lean~4 and Mathlib}
\label{sec:bg-lean}

\Lean is an interactive theorem prover and programming language based on
dependent type theory~\cite{demoura2021lean4}. Proofs and programs share the
same language, and correctness is guaranteed by a small trusted kernel that
independently type-checks every proof term.

The \Mathlib library~\cite{mathlib2020} provides formalized mathematics
including abstract algebra, analysis, and topology. For our development,
the key components are:
\begin{itemize}
  \item \texttt{NormedRing}, \texttt{NormedAlgebra}: typeclasses for
    normed rings and algebras with submultiplicative norms.
  \item \texttt{NormedSpace.exp}: the exponential function
    $a \mapsto \sum_{k} a^k/k!$ in a complete normed algebra.
  \item \texttt{intervalIntegral}: the Bochner integral over $[a,b]$,
    with the fundamental theorem of calculus.
  \item \texttt{HasDerivAt}: the Fr\'{e}chet derivative, with product rules
    and chain rules for normed algebra-valued functions.
\end{itemize}

\paragraph{Related formalization work.}
Formal verification has recently gained traction in mathematical physics and
quantum information.
Zhao and Yu~\cite{zhao2025chsh} formalized robust CHSH rigidity in \Lean,
discovering a gap in the original proof of McKague~(2012) during the
formalization process.
Meiburg et al.~\cite{meiburg2025steinlemma} formalized a generalized
quantum Stein's lemma, and
Tooby-Smith~\cite{toobysmith2026higgs} formalized the stability of the
two-Higgs-doublet potential, also identifying a literature error.
These projects share the pattern of using formal verification not merely to
certify known results but to expose subtle gaps in published proofs.
Our work extends this program to quantum simulation, providing the first
machine-checked proof of any product formula for operator exponentials.
